%!TEX root = ../template.tex
%% ====================================================================
%% Capitulo 3 -- Revisao sistematica do ambito
%% Esqueleto gerado 2026-08-13 a partir do indice v1
%% (Mind Palace: salas/tese/2026-08-13-indice-dissertacao-v1.md)
%% Cada '>>' descreve o que a seccao deve conter. Apagar ao escrever.
%% ====================================================================

\chapter{Revisão sistemática do âmbito}
\label{cha:revisao}

\section{Protocolo e pré-registo}
\label{sec:protocolo-sr}

%% >> Scoping review registada no OSF segundo PRISMA-ScR, com SQ1--SQ4 e a estrutura de dois corpora.


\section{Estratégia de pesquisa e seleção}
\label{sec:pesquisa}

%% >> Strings calibradas, screening em tres niveis, extracao assistida por dois modelos independentes sob protocolo auditado (Anexo ref app:validacao).


\section{Análise bibliométrica (Corpus A)}
\label{sec:corpus-a}

%% >> Distribuicao da literatura EnPI/EnB por ano, setor e contexto normativo: a evidencia quantitativa do gap.


\section{Síntese narrativa (Corpus B)}
\label{sec:corpus-b}

%% >> O que os estudos em processo continuo propoem, reconhecem e nao resolvem, por sub-questao.


\section{Síntese: o \emph{gap}}
\label{sec:gap}

%% >> Nao existe metodologia sistematica para avaliar e construir baselines energeticas em processo continuo dentro do enquadramento normativo.


